% ============================================
% MEM-16
% Gerado de molambudos.md — build_miolo.py
% ============================================
\section*{\textbf{MEM-16} --- A Morte no Colônia}
\markboth{MEM-16 --- A Morte no Colônia}{}

\noindent \textit{A última sessão}

\noindent \textbf{Junho de 1979 --- A última sessão}

\vspace{0.5em}\noindent\rule{\textwidth}{0.5pt}\vspace{0.5em}

\noindent Dr. Oliveira chegou na sala de sessões com o diário na mão. As folhas estavam marcadas. Algumas páginas estavam dobradas. Ele tinha lido tudo.

\noindent \textit{\textquotedblleft{}Joaquim\textquotedblright{}\textit{, ele disse. }\textquotedblleft{}Eu preciso entender uma coisa.\textquotedblright{}}

\noindent \textit{\textquotedblleft{}Pergunte.\textquotedblright{}}

\noindent \textit{\textquotedblleft{}A criatura. Ela está aqui agora?\textquotedblright{}}

\noindent \textit{\textquotedblleft{}Está.\textquotedblright{}}

\noindent \textit{\textquotedblleft{}Onde?\textquotedblright{}}

\noindent \textit{\textquotedblleft{}O senhor não vê ela. Mas ela vê o senhor.\textquotedblright{}}

\noindent Ele engoliu seco. Olhou para os lados.

\noindent \textit{\textquotedblleft{}O que ela está fazendo?\textquotedblright{}}

\noindent \textit{\textquotedblleft{}Ela está...\textquotedblright{}\textit{ Eu fechei os olhos. Senti o peso dentro de mim. O calor. A coisa se movendo entre meus órgãos. }\textquotedblleft{}Ela está se preparando.\textquotedblright{}}

\noindent \textit{\textquotedblleft{}Preparando para quê?\textquotedblright{}}

\noindent \textit{\textquotedblleft{}Para sair.\textquotedblright{}}

\noindent \textit{\textquotedblleft{}Sair? Sair como?\textquotedblright{}}

\noindent \textit{\textquotedblleft{}O senhor leu o diário. O senhor sabe como.\textquotedblright{}}

\noindent Dr. Oliveira ficou pálido.

\noindent \textit{\textquotedblleft{}Ela sai através de quem lê.\textquotedblright{}}

\noindent \textit{\textquotedblleft{}Isso.\textquotedblright{}}

\noindent \textit{\textquotedblleft{}Não.\textquotedblright{}}

\noindent \textit{\textquotedblleft{}Já é tarde, doutor. O senhor leu. Ela está dentro do senhor agora.\textquotedblright{}}

\noindent Ele balançou a cabeça.

\noindent \textit{\textquotedblleft{}Não. Não acredito nisso. Isso é sugestão. É o seu delírio. Você colocou essa ideia na minha cabeça porque...\textquotedblright{}}

\noindent \textit{\textquotedblleft{}Porque?\textquotedblright{}}

\noindent Ele não terminou.

\noindent A criatura começou a subir.

\noindent Eu senti ela saindo do buraco no meu estômago. Subindo pelo esôfago. Enchendo o peito. Os braços. As pernas. Os olhos.

\noindent Meu olho direito ardeu.

\noindent \textit{\textquotedblleft{}Doutor\textquotedblright{}\textit{, eu disse. A voz já não era toda minha. }\textquotedblleft{}Doutor, saia daqui.\textquotedblright{}}

\noindent \textit{\textquotedblleft{}O quê?\textquotedblright{}}

\noindent \textit{\textquotedblleft{}Saia. Agora.\textquotedblright{}}

\noindent Ele não saiu.

\noindent Ele ficou olhando para mim enquanto meu corpo começava a mudar. A pele do rosto esticou. Os dentes se apertaram. O olho amarelo começou a brilhar.

\noindent A criatura sorriu com a minha boca.

\noindent \textit{\textquotedblleft{}Você queria ver, doutor.\textquotedblright{}}

\noindent \textit{\textquotedblleft{}...\textquotedblright{}}

\noindent \textit{\textquotedblleft{}Então vê.\textquotedblright{}}

\vspace{0.5em}\noindent\rule{\textwidth}{0.5pt}\vspace{0.5em}

\noindent O que aconteceu depois eu não lembro direito.

\noindent Lembro de uma luz. Muito forte. Lembro do som de vidro quebrando. Lembro da cadeira caindo. Lembro de uma voz gritando --- era a minha voz? Era a voz dele?

\noindent Depois, o escuro.

\noindent Quando acordei, estava no chão. O diário estava aberto ao meu lado. O Dr. Oliveira estava encostado na parede, com os olhos arregalados, a boca aberta. Ele não se mexia.

\noindent Ele respirava.

\noindent Mas ele não estava mais ali.

\noindent Os olhos dele estavam vazios.

\noindent O médico que entrou na sala depois disse que eu tive uma convulsão. Que o Dr. Oliveira levou um susto. Que não foi nada.

\noindent Mas eu vi.\NE{Esta passagem descreve o mecanismo central de contaminação da narrativa: a criatura não se transfere por contato físico, mas por mediação textual. O leitor que acompanha o relato de Joaquim está, como Oliveira, exposto ao mesmo vetor de contágio. A tensão entre saber e não saber — entre ler e ser lido — é o conflito que o livro propõe a cada fragmento. O participante da leitura não é espectador: é cúmplice.}

\noindent A criatura tinha passado.

\noindent Ela estava saindo de mim.

\noindent E entrando nele.

\vspace{0.5em}\noindent\rule{\textwidth}{0.5pt}\vspace{0.5em}

\noindent Morri três dias depois.

\noindent Foi numa quarta-feira, 12 de junho de 1979. O coração parou. Disseram que foi parada cardíaca. Disseram que foi de velhice --- eu tinha 72 anos.

\noindent Eu sei o que foi.

\noindent A criatura saiu. Levou o que restava de mim.

\noindent O diário ficou.

\noindent Enterraram comigo. Ficou 62 dias na terra --- um dia para cada ano de prisão.

\noindent Quando me exumaram, o diário estava intacto.

\noindent As páginas ainda estavam úmidas.

\noindent A tinta ainda fresca.

\noindent \textit{\textquotedblleft{}O diário não acaba quando o autor morre\textquotedblright{}\textit{, eu escrevi na última página. }\textquotedblleft{}O diário começa quando o leitor abre.\textquotedblright{}}

\noindent Você abriu.

\noindent A criatura está lendo com você.

\vspace{0.5em}
\noindent\textit{\textbf{↪ Rotas:} \rota{DOC-06} \rotasep \rota{DOC-03} \rotasep \rota{LUC-12}}